\environment xi-style
\starttext
\title{拾穗者}

我时常梦见一片无边无际的田野。田里长的不是庄稼，是词语。它们有的饱满如麦穗，像“爱”“时间”“土地”这样沉甸甸的，弯着腰；有的干瘪细碎，像“的”“了”“在”，被风一吹就散落一地。我手里提着一只旧布囊，在田里慢慢地走，看见中意的词语，就弯腰拾起来，放进囊中。走完一整片田，我便在田埂上坐下来，把布囊里的东西倒出来点数，想知道这片田野到底盛产哪些词语。

这便是我写这个小程序的全部缘由。它没有任何高深的算法，只是一个在词语田野里低头走路的老人。它的正式名字叫“拾穗者”，其实不过是一个词频统计工具：你给它一篇文字，它还你一张词语的丰收清单。然而于我而言，那些代码片段就是拾穗的步骤，田野、囊、点数的动作，都实实在在地长在程序里。

\section{词语的穗子}

在田野里，每个词语都是一株独立的穗子。我需要一个容器来暂时存放它们——当然也就是一个链表，每个结点存着词语本身和它出现的次数。我们把它叫“穗子”，学名不妨叫 \type{WordNode}。

\startC
/BTEX\color[darkred]{@}/ETEX/BTEX\color[darkred]{ 穗子的形状 }/ETEX/BTEX\color[darkred]{\#}/ETEX/BTEX\reference[xi-1]{1}\inoutermargin{\darkred{1}}/ETEX
typedef struct word_node {
        char *word;           /* 词语本身 */
        int count;            /* 被拾起的次数 */
        struct word_node *next;
} WordNode;
/BTEX=>\color[darkyellow]{ 主程序：拾穗的过程 }<\in[xi-5]>/ETEX
\stopC

\noindent 田野广阔，我不会一直盯着同一株穗子。每遇到一个词语，我就得在布囊里翻找一下，看看是不是已经拾起过。布囊就是整个链表，我们用 \type{find_word} 这个动作来完成翻找。

\startC
/BTEX\color[darkred]{@}/ETEX/BTEX\color[darkred]{ 在布囊中翻找 }/ETEX/BTEX\color[darkred]{\#}/ETEX/BTEX\reference[xi-2]{2}\inoutermargin{\darkred{2}}/ETEX
WordNode *find_word(WordNode *head, const char *word) {
        WordNode *p = head;
        while (p) {
                if (strcmp(p->word, word) == 0)
                        return p;
                p = p->next;
        }
        return NULL;
}
/BTEX=>\color[darkyellow]{ 主程序：拾穗的过程 }<\in[xi-5]>/ETEX
\stopC

\noindent 如果布囊里还没有这株穗子，我就得新添一个结点，把它放进去。这个动作叫 \type{add_word}，它总是把新穗子放在链表的开头，因为刚拾起的总是最新鲜的。

\startC
/BTEX\color[darkred]{@}/ETEX/BTEX\color[darkred]{ 添加新的穗子 }/ETEX/BTEX\color[darkred]{\#}/ETEX/BTEX\reference[xi-3]{3}\inoutermargin{\darkred{3}}/ETEX
WordNode *add_word(WordNode *head, const char *word) {
        WordNode *node = malloc(sizeof(WordNode));
        node->word = strdup(word);
        node->count = 1;
        node->next = head;
        return node;
}
/BTEX=>\color[darkyellow]{ 主程序：拾穗的过程 }<\in[xi-5]>/ETEX
\stopC

\section{分割词语的落叶}

田野里的文字不是天然就分成一株一株穗子的。它是一片连绵的文本，我需要一个分割的动作，把词语从句子中分离出来。这个动作不复杂：我只关心字母和数字，其余的一概视为分隔符。这一步，仿佛是用一把竹耙将落叶聚拢，只留下有用的叶片。\type{get_word} 函数每次从输入里扒出一个词，填进我递过去的小竹篮。

\startC
/BTEX\color[darkred]{@}/ETEX/BTEX\color[darkred]{ 扒出一个词 }/ETEX/BTEX\color[darkred]{\#}/ETEX/BTEX\reference[xi-4]{4}\inoutermargin{\darkred{4}}/ETEX
int get_word(FILE *fp, char *buf, int buf_size) {
        int c;
        /* 跳过非字母数字字符 */
        while ((c = fgetc(fp)) != EOF) {
                if (isalnum(c))
                        break;
        }
        if (c == EOF) return 0;

        int i = 0;
        buf[i++] = tolower(c);
        while (i < buf_size - 1 && (c = fgetc(fp)) != EOF && isalnum(c)) {
                buf[i++] = tolower(c);
        }
        buf[i] = '\0';
        return 1;
}
/BTEX=>\color[darkyellow]{ 主程序：拾穗的过程 }<\in[xi-5]>/ETEX
\stopC

\section{走在田野里}

现在，一切准备就绪。我打开田野的大门——也许就是一个文件，也许是标准输入，那便是另一片更广阔的田地。我拎着空空的布囊，一步一步走进去，每遇到一个词，就拾起来放进囊中。走完全部田地，我便坐在田埂上清点。

主程序逻辑极其简单，像一段默默行路的时光。它所需的外部世界（那些头文件）早已在“外景”里备好，此刻只需轻轻引出。

\startC
/BTEX\color[darkred]{@}/ETEX/BTEX\color[darkred]{ 主程序：拾穗的过程 }/ETEX/BTEX\color[darkred]{\#}/ETEX/BTEX\reference[xi-5]{5}\inoutermargin{\darkred{5}}/ETEX
/BTEX\color[darkcyan]{\#}/ETEX/BTEX\color[darkcyan]{ 外景 }/ETEX/BTEX\color[darkcyan]{@}/ETEX/BTEX\ <\in[xi-6]>/ETEX
/BTEX\color[darkcyan]{\#}/ETEX/BTEX\color[darkcyan]{ 穗子的形状 }/ETEX/BTEX\color[darkcyan]{@}/ETEX/BTEX\ <\in[xi-1]>/ETEX
/BTEX\color[darkcyan]{\#}/ETEX/BTEX\color[darkcyan]{ 在布囊中翻找 }/ETEX/BTEX\color[darkcyan]{@}/ETEX/BTEX\ <\in[xi-2]>/ETEX
/BTEX\color[darkcyan]{\#}/ETEX/BTEX\color[darkcyan]{ 添加新的穗子 }/ETEX/BTEX\color[darkcyan]{@}/ETEX/BTEX\ <\in[xi-3]>/ETEX
/BTEX\color[darkcyan]{\#}/ETEX/BTEX\color[darkcyan]{ 扒出一个词 }/ETEX/BTEX\color[darkcyan]{@}/ETEX/BTEX\ <\in[xi-4]>/ETEX

int main(int argc, char *argv[]) {
        FILE *fp = stdin;
        if (argc > 1) {
                fp = fopen(argv[1], "r");
                if (!fp) {
                        fprintf(stderr, "无法打开田野 %s\n", argv[1]);
                        return 1;
                }
        }

        WordNode *head = NULL;
        char buf[256];

        while (get_word(fp, buf, sizeof(buf))) {
                WordNode *node = find_word(head, buf);
                if (node) {
                        node->count++;
                } else {
                        head = add_word(head, buf);
                }
        }

        if (fp != stdin) fclose(fp);

        /* 清点：找出出现次数最多的若干词语 */
        int top_n = 20;
        WordNode *p;
        for (int i = 0; i < top_n && head; i++) {
                WordNode *max_node = head;
                for (p = head->next; p; p = p->next) {
                        if (p->count > max_node->count)
                                max_node = p;
                }
                if (max_node->count == 0) break;
                printf("%s: %d\n", max_node->word, max_node->count);
                max_node->count = 0;
        }

        /* 释放布囊，离开田野 */
        while (head) {
                WordNode *next = head->next;
                free(head->word);
                free(head);
                head = next;
        }

        return 0;
}
\stopC

\noindent 你也许注意到了，我的清点方式很笨拙：每次循环都遍历整个链表，找出当前最高的一个，然后把它标记为零，再找下一个。这就像在田埂上把穗子一遍遍摊开，每次只挑出最饱满的那一株。我完全可以用更精巧的方法，但一个老人坐在田埂上，慢慢地挑拣，本身也是拾穗的一部分。程序只要能用，便不必总是聪明绝顶。

\section{外景：必须包含的头文件}

每个程序都得有一些外来的支撑，就像我脚下这片田野，总得承认天空和远方山脉的存在。这些头文件正是我的田野之外的外景。

\startC
/BTEX\color[darkred]{@}/ETEX/BTEX\color[darkred]{ 外景 }/ETEX/BTEX\color[darkred]{\#}/ETEX/BTEX\reference[xi-6]{6}\inoutermargin{\darkred{6}}/ETEX
#include <stdio.h>
#include <stdlib.h>
#include <string.h>
#include <ctype.h>
/BTEX=>\color[darkyellow]{ 主程序：拾穗的过程 }<\in[xi-5]>/ETEX
\stopC

\noindent 它们的名字连在一起，看起来像一扇窄窄的栅栏门。推开它，词语田野就在眼前了。

\section{运行拾穗者}

把上面的代码片段提取出来，编译，就得到了一个叫 \type{gleaner} 的小东西。你可以这样用它：

\starttyping[color=darkcyan]
$ gleaner my_text.txt
\stoptyping

它会输出文章里出现最多的二十个词，就像老人把布囊里的穗子摊开给你看。如果什么文件也不给，它就安静地等着你从键盘上播种——等你敲完，按下 \type{Ctrl+D}，它便照样清点你刚才亲手撒下的词语。

对我来说，这便是一种诚实的文学编程了：叙事和代码谁也不为谁服务，它们本就是同一种劳作的两个面向——拾穗的动作与拾穗的故事。代码片段就是那些穗子，而文字则是拾穗人晚归时哼唱的歌。
\stoptext